\documentclass{report}
\usepackage{amsmath,amssymb}
\setlength{\parindent}{0mm}
\setlength{\parskip}{1em}
\begin{document}
\begin{center}
\rule{6in}{1pt} \
{\large Pieter Ghysels \\
{\bf Increasing the Arithmetic Intensity of Multigrid on Many-Core Chips}}

Middelheimcampus M G 320b \\ Middelheimlaan 1 \\ 2020 Antwerpen
\\
{\tt pieter.ghysels@ua.ac.be}\\
Wim I. Vanroose\end{center}

The basic building blocks of a classic multigrid algorithm, which are
essentially stencil computations, all have a low ratio of executed
floating point operations per byte fetched from memory. This important
ratio can be identified as the arithmetic intensity. According to the
roofline model [1], applications with a low arithmetic intensity are
typically bounded by memory traffic, whereas applications with higher
arithmetic intensity are bounded by the floating point unit throughput.
It is well known that many bandwidth limited algorithms typically perform
at a small percentage of the theoretical peak performance of the
underlying hardware.

In this work we look at changes to both the multigrid algorithm and its
implementation in order to increase the arithmetic intensity.

We present a theoretical model for the cost of a multigrid cycle (V-cycle
or full multigrid) that includes the cost of communication, based on the
roofline model. Our multigrid model predicts performance benefits for
code optimization strategies such as: fusion of operations, tiling,
parallelization, vectorization (SIMDization) and cache prefetching. The
multigrid convergence rate is taken into account in the model for
different smoothers, different coarsening and interpolation strategies,
different cycles and different coarsest grid solvers.

Our theoretical model is validated using numerical experiments on modern
multi and many-core hardware. Instead of resorting to manual and tedious
low level code optimizations, we piggyback on recent compiler
improvements, such as stencil compilers, automatic
parallelizers and automatic loop transformers for cache locality
optimization [2,3]. An important observation is that the underlying
hardware can only be used efficiently when the number of consecutive
smoothing steps is large enough. In that case, the arithmetic intensity
can be increased sufficiently such that SIMD vector units can be kept
busy and the extra smoothing steps pay off in the overall performance of
the algorithm. It is thus important to look for algorithmic changes that
can put these extra smoothing steps to
effective use.

[1] Williams S, Waterman A and Patterson DA. Roofline: an insightful
Visual Performance model for multicore architectures. Communications of
the ACM 2009; 52(4):65�76

[2] Christen M, Schenk O, and Burkhart H. Patus: A code generation and
autotuning framework for parallel iterative stencil computations on
modern microarchitectures. In Parallel \& Distributed Processing
Symposium (IPDPS), 2011 IEEE International, pages 676�687. IEEE, 2011.

[3] Bondhugula U, Baskaran M, Krishnamoorthy S, Ramanujam J, Rountev A,
and Sadayappan P. Automatic Transformations for Communication-Minimized
Parallelization and Locality Optimization in the Polyhedral Model.
International Conference on Compiler Construction (ETAPS CC), Apr 2008,
Budapest, Hungary.


\end{document}
